\chapter{Lernziele}

\begin{todolist}
    \item Ich kann die Grundkonzepte und Unterschiede zwischen IP, IPv4 und IPv6 erklären.
    \item Ich kann für IPv4 und IPv6 jeweils die Anzahl möglicher Adressen angeben und berechnen.
    \item Ich kann die Funktionen und Arbeitsweisen von Routern und Switches in einem Netzwerk verstehen und erläutern.
    % \item Ich kann eine Routing-Tabelle interpretieren und für einfache Routing-Entscheidungen nutzen.
    % \item Ich kann eine Subnetzmaske erstellen und anwenden, um Netzwerke in verschiedene Subnetze zu unterteilen.
    \item Ich kann die grundlegenden Konzepte von TCP beschreiben und dessen Einsatz im Kontext von Netzwerkprotokollen verstehen.
    % \item Ich kann grundlegende HTML-Strukturen erstellen und verstehen, wie HTML-Seiten über HTTP übertragen werden.
    \item Ich kann den Ablauf eines Abrufs von Webinhalten über das Internet in einzelnen Schritten und den darin involvierten Protokollen erläutern.
    \item Ich kann die Rolle eines DNS-Servers im Netzwerk erklären und wie DNS-Anfragen bearbeitet und beantwortet werden.
    \item Ich kann grundlegende Netzwerkbefehle wie \texttt{ping}, \texttt{ipconfig}/\texttt{ifconfig} und \texttt{nslookup} anwenden und deren Ausgabe interpretieren.
    \item Ich kann einfache Netzwerkprobleme identifizieren und grundlegende Schritte zur Fehlerbehebung durchführen.
    \item Ich kann die Client-Server-Architektur erklären und deren Bedeutung für die Kommunikation im Internet beschreiben.
    \item Ich kann in Python einfache Socket-Programme schreiben, die eine Verbindung zu einem Server herstellen und Nachrichten senden/empfangen.
    \item Ich kann eine Chat-Anwendung in Python entwickeln, die über einen Relay-Server Nachrichten zwischen mehreren Clients austauscht.
    \item Ich kann die grundlegenden Schritte zur Fehlerbehebung einer nicht funktionierenden Internet-Verbindung durchführen und dokumentieren.
\end{todolist}
